
% IEEE conference version for GTSD 2026 proceedings
\documentclass[conference]{IEEEtran}
\usepackage{cite}
\usepackage{amsmath,amssymb,amsfonts}
\usepackage{graphicx}
\graphicspath{{figures/}}
\usepackage{textcomp}
\def\BibTeX{{\rm B\kern-.05em{\sc i\kern-.025em b}\kern-.08em
    T\kern-.1667em\lower.7ex\hbox{E}\kern-.125emX}}

\begin{document}
\bstctlcite{IEEEexample:BSTcontrol}

\title{Learning-Based Indoor Visible Light Positioning for Estimating 3D Objects}

\author{\IEEEauthorblockN{Quoc-Phong Pham, Le The Dung}
\IEEEauthorblockA{\textit{Faculty of Electrical and Electronics Engineering}\\
\textit{Ho Chi Minh City University of Technology and Engineering}\\
Ho Chi Minh City, Viet Nam\\
2530713@student.hcmute.edu.vn, dunglt@hcmute.edu.vn}
}
\maketitle

\begin{abstract}
In this paper, we propose an integrated data-driven framework to enhance the precision of visible light positioning (VLP) systems under complex occlusion and noise conditions. Specifically, we utilize a direct current line-of-sight (LoS) optical channel model combined with a geometric occlusion mechanism for generating received signal strength (RSS) maps across a photodiode (PD) sensor grid. Then, we employ a multi-layer perceptron (MLP) network to perform multi-output regression, learning the non-linear mapping from logarithmically transformed RSS features to target parameters, including planar position, radius, and height. The efficacy of the model is validated through root mean square error (RMSE) analysis, cumulative distribution function (CDF) of errors, trajectory tracking, and noise robustness stress tests. The proposed model achieves a test planar localization RMSE of 6.75 cm, corresponding to approximately 1.01 percent of the positioning-area diagonal. The RMSE values for radius and height estimation are 1.14 cm and 8.44 cm, respectively. The results indicate that blockage-aware RSS maps can provide informative spatial signatures for object-aware VLP, although height estimation remains more challenging because of weaker vertical observability. These findings establish a promising learning-based approach for mitigating occlusion-induced errors and supporting high-precision indoor object estimation in VLP environments.
\end{abstract}

\begin{IEEEkeywords}
Multi-Layer Perceptron (MLP), Visible Light Positioning (VLP), indoor positioning.
\end{IEEEkeywords}




\section{Proposed Method}

The proposed method estimates the object parameter vector $\mathbf{y}=[x,y,r,h]$ from blockage-affected RSS observations over the PD grid. A simulation-driven dataset is first generated by randomly sampling valid object parameters inside the room. For each sample, the LoS channel model and the blockage condition described in Section~II are evaluated over all PD locations for all four LEDs, producing four RSS maps that encode the spatial attenuation pattern caused by the object. Gaussian noise is added to the received power before transforming the measurements into the logarithmic RSS domain.

The overall processing pipeline is summarized in Fig.~\ref{fig:pipeline}. Starting from the room geometry and sampled object configuration, four LED-specific RSS maps are stacked into a structured tensor, flattened into a feature vector, standardized, and finally processed by a multi-output MLP regressor to estimate $[x,y,r,h]$.

\begin{figure}[tb]
\centering
\includegraphics[width=\columnwidth]{figure_16.png}
\caption{Overview of the proposed framework: scene sampling, RSS tensor construction, feature flattening and scaling, MLP regression, and estimation of $[x,y,r,h]$.}
\label{fig:pipeline}
\end{figure}

The blockage effect is visualized in Fig.~\ref{fig:rss_heatmaps_2d}, where each LED produces a different attenuation pattern depending on the relative transmitter--object--receiver geometry. These complementary shadow patterns are the main features exploited by the regressor.

\begin{figure}[tb]
\centering
\includegraphics[width=\columnwidth]{figure_05_rss_heatmaps_2d.png}
\caption{Top-view RSS maps for the four LEDs. Dark regions indicate severe attenuation caused by object blockage.}
\label{fig:rss_heatmaps_2d}
\end{figure}

For each object realization, the RSS maps generated by the four LEDs are concatenated into a single feature vector. Measurement uncertainty is modeled by adding Gaussian noise before logarithmic transformation:
\begin{equation}
P_r^{\mathrm{noisy}}=\max\!\left(P_r+\eta,\varepsilon\right),\qquad \eta\sim\mathcal{N}(0,\sigma^2),
\end{equation}
followed by
\begin{equation}
P_{\mathrm{RSS}}=10\log_{10}\!\left(P_r^{\mathrm{noisy}}\right).
\end{equation}
The learning task is therefore formulated as a nonlinear regression problem,
\begin{equation}
f:\mathbf{x}\mapsto \mathbf{y},\qquad \mathbf{x}\in\mathbb{R}^{N_f},\;\mathbf{y}=[x,y,r,h]\in\mathbb{R}^{4},
\end{equation}
where $\mathbf{x}$ is the concatenated RSS feature vector and $\mathbf{y}$ is the target parameter vector.

The RSS maps from the four LEDs are then concatenated into a single input vector,
\begin{equation}
\mathbf{x}=[\mathbf{x}_1,\mathbf{x}_2,\mathbf{x}_3,\mathbf{x}_4]\in\mathbb{R}^{N_f},
\end{equation}
while the corresponding target output is
\begin{equation}
\mathbf{y}=[x,y,r,h]\in\mathbb{R}^{4}.
\end{equation}
Before training, both inputs and targets are normalized as
\begin{equation}
\tilde{\mathbf{x}}^{(i)}=\frac{\mathbf{x}^{(i)}-\boldsymbol{\mu}_x}{\boldsymbol{\sigma}_x},
\qquad
\tilde{\mathbf{y}}^{(i)}=\frac{\mathbf{y}^{(i)}-\boldsymbol{\mu}_y}{\boldsymbol{\sigma}_y},
\end{equation}
which improves numerical stability and balances the scales of the regression targets.

To learn the nonlinear mapping between RSS patterns and object parameters, we employ a multilayer perceptron (MLP) with four hidden layers containing 512, 256, 128, and 64 neurons, respectively, using ReLU activation. Let $\tilde{\mathbf{x}}$ be the normalized input vector. The hidden layers are computed recursively as
\begin{equation}
\mathbf{z}^{(l)}=\sigma\!\left(\mathbf{W}^{(l)}\mathbf{z}^{(l-1)}+\mathbf{b}^{(l)}\right),\qquad l=1,\dots,L-1,
\end{equation}
with $\mathbf{z}^{(0)}=\tilde{\mathbf{x}}$. The output layer is linear,
\begin{equation}
\hat{\mathbf{y}}=\mathbf{W}^{(L)}\mathbf{z}^{(L-1)}+\mathbf{b}^{(L)},
\end{equation}
so that the model directly predicts continuous-valued object parameters.

Training is performed by minimizing the mean squared error (MSE)
\begin{equation}
\mathcal{L}_{\mathrm{MSE}}=\frac{1}{N}\sum_{i=1}^{N}\left\|\hat{\mathbf{y}}^{(i)}-\tilde{\mathbf{y}}^{(i)}\right\|_2^2,
\end{equation}
using Adam optimization with mini-batches and early stopping. During inference, a test sample is normalized using the statistics of the training set, passed through the trained MLP, and then inverse-transformed to obtain the final estimate
\begin{equation}
\mathbf{y}_{\mathrm{pred}}=\hat{\mathbf{y}}\odot\boldsymbol{\sigma}_y+\boldsymbol{\mu}_y.
\end{equation}
This yields the predicted planar position and geometric parameters of the object in a single forward pass.

\section{Experimental Setup and Performance Evaluation}

This section evaluates the proposed learning-based indoor VLP framework under blockage and noise conditions.

\subsection{Experimental Setup}

The simulation environment follows the configuration described in Section~II. The room size is $5\times 5\times 3~\mathrm{m}^3$, the LED semi-angle at half power is $60^\circ$, the photodiode active area is $10^{-4}~\mathrm{m}^2$, the receiver field of view is $60^\circ$, and the optical concentrator refractive index is 1.5. The PD grid resolution is $0.2~\mathrm{m}$. Object parameters are sampled uniformly within
\begin{equation}
r\in[0.15,0.40]~\mathrm{m},\qquad h\in[1.2,1.9]~\mathrm{m},
\end{equation}
and the dataset contains $50{,}000$ simulated samples split into training, validation, and test sets with a ratio of $80\%/10\%/10\%$. Gaussian noise with standard deviation $10^{-8}$ is added to the received optical power before the logarithmic RSS transformation. Table~\ref{tab:simulation_parameters} summarizes the main simulation and learning parameters used in the evaluation.


\IEEEtriggeratref{7}
\bibliographystyle{IEEEtran}
\bibliography{bib/references}

\end{document}
